\section{The bug a unit test cannot see}
\label{sec:interaction}

A \textbf{unit test} tests one unit in isolation. A \textbf{component test} tests several parts
working together, still without hardware: here \code{app::EchoNode} plus
\code{driver::can::Stub}, connected through \code{driver::can::Interface}.

The difference shows most clearly in the kind of bug they catch. Suppose \code{run()} is written
like this:

\begin{cppcode}
Frame frame{};
if (!myDriver.receive(frame)) { return; }

Frame response{};
response.id  = myResponseId;
response.dlc = myDriver.receive(frame) ? frame.dlc : 0U;   // wrong: calls receive() again
\end{cppcode}

Every unit test is still green. \code{Stub} does exactly what it promises, and so does
\code{Frame}. The bug is in the \textbf{interaction}: the second call consumes the state the first
one left behind.

\textbf{That is the whole point.} Interaction bugs --- one call too many, a call in the wrong
order, an ignored return value --- are the most common kind of fault in driver code, and they live
between the units, not in them.

\section{\texorpdfstring{\code{app::EchoNode}}{app::EchoNode}}
\label{sec:echonode}

An application component that does one thing: reads a frame and echoes it back with a different
ID.

\begin{cppcode}
namespace driver::can { class Interface; }

namespace app
{
class EchoNode
{
public:
    explicit EchoNode(driver::can::Interface& driver, std::uint16_t responseId) noexcept;

    void run() noexcept;
    [[nodiscard]] std::uint32_t echoCount() const noexcept;

private:
    driver::can::Interface& myDriver; /**< Driver used for all traffic. Not owned. */
    std::uint16_t myResponseId;
    std::uint32_t myEchoCount;
};
} // namespace app
\end{cppcode}

\subsection{The layer boundary, in one line}

Note what the header does \emph{not} include: \code{stub.hpp}, \code{spi.hpp},
\code{byte_transport.hpp}. Only a forward declaration.

\textbf{That line is the architecture's boundary expressed in code.} A reviewer who sees
\code{#include "driver/can/spi.hpp"} here should stop the PR on the spot. The rule, which holds
for the whole project: no file above \code{driver/can/interface.hpp} may contain the word SPI.

\subsection{Dependency injection}

\code{EchoNode} \textbf{does not create} its driver. It is given one. The alternative is worth
looking at to see what it costs:

\begin{cppcode}
class EchoNode
{
public:
    EchoNode() : myDriver{} {}   // creates its own Spi
private:
    driver::can::Spi myDriver;   // no
};
\end{cppcode}

That variant is impossible to test: there is no way to give it a stub, and therefore no way to run
it without an FPGA on the desk. It also binds the application to a concrete driver type, which
makes the whole \code{Interface} pointless.

\textbf{Why a reference and not a pointer?} A reference cannot be null and cannot be swapped after
construction. Both are properties we want.

\textbf{Why not \code{std::unique_ptr}?} Because \code{EchoNode} does not \emph{own} the driver.
The ownership sits with the factory in \kapref{ch:avr}, and the driver's lifetime is longer than
the component's. The Doxygen comment explicitly says "Not owned", and that is worth writing out.

\subsection{Two details in \texorpdfstring{\code{run()}}{run()}}

\textbf{The loop runs to \code{frame.dlc}, not to \code{MaxDataLen}.} Copying eight bytes would
have "worked" since the rest are zeroes anyway --- but \code{response.dlc} says how many apply,
and code that copies more than it claims to copy stops being right when some other assumption
changes. Exactly the same reasoning as the masking in \kapref{ch:driver}.

\textbf{\code{myEchoCount} is incremented only when \code{send()} succeeded.} The counter means
"frames echoed", not "attempts".

\section{Reading the component test}
\label{sec:readcomp}

\begin{cppcode}
TEST(EchoNode, echoesInjectedFrameWithResponseId)
{
    // Arrange.
    driver::can::Stub stub{};
    app::EchoNode node{stub, 0x200U};

    driver::can::Frame incoming{};
    incoming.id      = 0x123U;
    incoming.dlc     = 2U;
    incoming.data[0] = 0xAAU;
    incoming.data[1] = 0xBBU;
    stub.inject(incoming);

    // Act.
    node.run();

    // Assert.
    driver::can::Frame echoed{};
    EXPECT_TRUE(stub.receive(echoed));
    EXPECT_EQ(echoed.id, 0x200U);
    EXPECT_EQ(node.echoCount(), 1U);
}
\end{cppcode}

The test has two roles at once: it \textbf{feeds in} traffic with \code{inject()}, and it
\textbf{inspects} outgoing traffic with \code{receive()} --- both through the same test double.

What the supplied suite covers:

\begin{booktable}{@{}lL@{}}
\thead{Case} & \thead{What it proves} \\ \midrule
No injected frame & \code{echoCount()} unchanged. \\
One injected frame & The right \code{id}, \code{dlc} and data are echoed; the counter goes up by
  one. \\
Two frames, two \code{run()} calls & The counter becomes 2. \\
One frame, \code{run()} \textbf{twice} & The counter becomes 1 --- not 2. Catches the bug in
  \secref{sec:interaction}. \\
\code{dlc = 0} and \code{dlc = 8} & A pair of boundary values. \\
\code{send()} fails & The counter does not go up. \\
\end{booktable}

\subsection{The same pattern, one layer down}

This is the only thing you need to take with you from here. In \kapref{ch:testning} the supplied
suite does \textbf{exactly the same thing}, one seam further down: \code{BankTransport} replaces
\code{Stub}, \code{Spi} is what is being tested, and the test scripts what the hardware answers
instead of which frames come in.

That it \emph{is} the same pattern is no coincidence --- it is what you get from putting the seams
in the right places.

\subsection{What it does not prove}

That \code{EchoNode} interacts correctly with \emph{some} implementation of \code{Interface}.
Nothing about \code{Spi}, which does not exist yet, and nothing about whether the register map has
been read correctly.
